\section{Setting and definitions}\label{sec:definitions}

\begin{definition}[Graph and adjacency; textbook Definitions 1.1.1 and 1.1.4]\label{def:graph}
A \emph{graph} is a pair $G = (V, E)$ of a non-empty finite set $V$ of \emph{vertices} and a set $E$ of \emph{edges}; it is \emph{simple undirected} when each edge is an unordered pair of two distinct vertices, so there are no loops and no repeated edges \cite[Definition 1.1.1, p.~9; Definition 1.1.4, p.~11]{petrosyan_graph_theory} \cite[Definition 1.1.1]{galway_adjacency_notes}. All graphs in this paper are finite, simple, and undirected. Write $V = \{n_1, \dots, n_N\}$ with $N \ge 1$ (the textbook writes $v_i$), and write $n_i n_j = n_j n_i$ for an edge. Two vertices are \emph{adjacent} if they form an edge.
\end{definition}

\begin{definition}[Adjacency matrix; textbook, p.~11, no number of its own]\label{def:matrix}
The \emph{adjacency matrix} $A$ of $G$ is the $N \times N$ matrix with $A_{ij} = 1$ if $n_i$ and $n_j$ are adjacent and $A_{ij} = 0$ otherwise \cite[p.~11]{petrosyan_graph_theory}. Edges are unordered pairs of distinct vertices, so $A_{ij} = A_{ji}$ and $A_{ii} = 0$ \cite[p.~11]{petrosyan_graph_theory}.
\end{definition}

For $N \times N$ matrices $X, Y$ the product has entries $(XY)_{ij} = \sum_{r=1}^{N} X_{ir} Y_{rj}$. The powers of $A$ are $A^0 = I$, the identity matrix (ones on the diagonal, zeros elsewhere), and $A^{k+1} = A^k A$. The entry in row $i$ and column $j$ of $A^k$ is written $(A^k)_{ij}$.

\begin{definition}[Walk and length; textbook Definition 1.2.10]\label{def:walk}
A \emph{walk} of \emph{length} $k \ge 0$ from $n_i$ to $n_j$ is a sequence $(v_0, v_1, \dots, v_k)$ of vertices with $v_0 = n_i$, $v_k = n_j$, and $v_{s-1}$ adjacent to $v_s$ for $1 \le s \le k$. The length $k$ is the number of edges counted with repetition.
\end{definition}

The textbook writes a walk as an alternating sequence of vertices and edges $u_0$, $e_1$, $u_1$, \dots, $e_k$, $u_k$, with $e_s = u_{s-1}u_s$ \cite[Definition 1.2.10, p.~21]{petrosyan_graph_theory}. In a simple graph two vertices carry at most one edge, so the edges are determined by the vertices and the two forms describe the same walks (a derived remark); Grinberg also writes walks as vertex sequences \cite[Definition 2.9.1(a), p.~36]{grinberg_graph_theory}. In a walk both vertices and edges may repeat \cite[p.~22]{petrosyan_graph_theory}. This paper allows $k = 0$: the single vertex $(n_i)$ is a walk of length $0$ from $n_i$ to itself. That is the paper's convention; it agrees with the textbook's property $d(u,v) = 0$ exactly when $u = v$ \cite[p.~23]{petrosyan_graph_theory} and with Grinberg's definition \cite[Example 2.9.2, p.~37]{grinberg_graph_theory}. Since $A$ is symmetric, a walk read backwards is a walk.

\begin{definition}[Path; textbook Definition 1.2.13, ``simple path'']\label{def:path}
A \emph{path} is a walk whose vertices are pairwise distinct. A \emph{path between $n_i$ and $n_j$}, written from $n_i$ to $n_j$, is a path $(v_0, \dots, v_k)$ with $v_0 = n_i$ and $v_k = n_j$ \cite[Definition 1.2.13, p.~22]{petrosyan_graph_theory} \cite[Definition 2.9.1(c),(d)]{grinberg_graph_theory} \cite{bai_bimsa_adjacency}. This is the meaning of ``path'' throughout the paper.
\end{definition}

\begin{remark}\label{rem:trap}
The textbook's own word ``path'' means something weaker: a walk whose edges are pairwise distinct \cite[Definition 1.2.12, p.~22]{petrosyan_graph_theory}, called a \emph{trail} in English. This paper never uses ``path'' in that sense; it writes ``trail''. By \Cref{def:path} each vertex $n_i$ has exactly one path to itself, $(n_i)$, of length $0$. Reversing a path from $n_i$ to $n_j$ gives a path from $n_j$ to $n_i$, so ``between'' is symmetric.
\end{remark}

\begin{definition}[Connected; textbook Definition 2.1.1, restated]\label{def:connected}
$G$ is \emph{connected} if for every pair $n_i, n_j$ of vertices there is a path from $n_i$ to $n_j$ \cite[Definition 2.1.1, p.~48]{petrosyan_graph_theory}.
\end{definition}

The textbook states this with its trail in place of the path; \Cref{prop:same} shows that the two statements describe the same graphs. For $i = j$ the length-$0$ path settles the condition, so only pairs with $i \ne j$ carry content.
